% Generated draft: requires mathematical and source review.
\documentclass[11pt]{article}
\usepackage[margin=25mm]{geometry}
\usepackage{fontspec}
\setmainfont{Latin Modern Roman}
\usepackage{amsmath,amssymb,mathtools,unicode-math}
\setmathfont{Latin Modern Math}
\usepackage{enumitem,needspace,xurl,hyperref}
\setlength{\parindent}{0pt}
\setlength{\parskip}{5pt}
\setlength{\emergencystretch}{3em}
\newcommand{\N}{\mathbb{N}}
\newcommand{\Z}{\mathbb{Z}}
\newcommand{\Q}{\mathbb{Q}}
\newcommand{\R}{\mathbb{R}}
\newcommand{\C}{\mathbb{C}}
\newcommand{\F}{\mathbb{F}}
\newcommand{\A}{\mathbb{A}}
\newcommand{\PP}{\mathbb{P}}
\newcommand{\E}{\mathbb{E}}
\newcommand{\eps}{\varepsilon}
\newcommand{\dd}{\,\mathrm{d}}
\begin{document}
\section*{UnsolvedMath ID: 7200023}
\textbf{Source problem code:} catalogue entry
\subsection{Definitions and mathematical statement}
Let \(k\) be a field of characteristic \(p>0\), and let \(X\) be an integral separated scheme of finite type over \(k\), of arbitrary finite dimension. In the regular-resolution interpretation, a resolution is a proper birational morphism
\[
 \pi:Y\to X
\]
from a regular integral variety \(Y\). Regular means that every local ring is regular; birational means that \(\pi\) induces an isomorphism of function fields. The problem asks whether every such \(X\) admits a resolution, for all positive characteristics and dimensions.

The source does not delimit the base fields or the strength of resolution. Over a perfect field regularity gives smoothness for finite-type varieties, but over an imperfect field a regular variety need not be smooth over \(k\). Thus replacing regular by smooth for arbitrary fields changes the target. Nor does the excerpt specify an isomorphism over the regular locus, a sequence of permissible blowups, embedded/logarithmic resolution, functoriality, or an algorithm. These stronger variants are not silently added. Alterations that are only generically finite are not birational resolutions of the displayed type.

The supplied summary reports characteristic-zero resolution attributed to Hironaka, positive-characteristic progress through dimension three, and the general dimension-at-least-four problem open as of its 2026 assessment. It attributes earlier threefold cases in characteristic greater than five to Abhyankar, later cases to Cossart--Piltant, and logarithmic/embedded work to Cossart--Jannsen--Saito. It gives no precise theorem scopes or full citations; its claims about all quasi-excellent schemes and embedded/logarithmic threefold resolution require verification and must not be equated automatically with regular variety resolution. The mentioned monomial-case and dimension-two work is special-case progress, not an arbitrary-dimensional solution.

The primary-source status audit remains outstanding in the supplied provenance. The pinned Wikipedia label and the later summary's report of no 2024--2026 general announcement are source assessments, not an independently verified current literature review.
\subsection{Short English statement}
Resolve every positive-characteristic algebraic variety in arbitrary dimension by a proper birational map from a regular variety. The field and resolution conventions matter, especially over imperfect fields; the supplied low-dimensional and stronger embedded-resolution claims require precise source checking.
\subsection{Sources}
\begin{itemize}
\item Local registry: \url{data/unsolvedmath/problems.json}, permanent UnsolvedMath ID \(7200023\), source code catalogue entry; supplied statement, background, and research summary, updated 2026-08-06.
\item Wikipedia contributors, \emph{List of unsolved problems in mathematics}, Geometry bullet 23, pinned revision 1366636767: \url{https://en.wikipedia.org/w/index.php?title=List_of_unsolved_problems_in_mathematics&oldid=1366636767#Geometry}. MediaWiki-source extraction accessed 2026-07-29; primary-source status review remains required.
\item Hironaka; Abhyankar; Cossart--Piltant; and Cossart--Jannsen--Saito are named in the summary without full citations or exact theorem hypotheses.
\item arXiv:1507.05195, mentioned for monomial-case progress with only its abstract reported seen; arXiv:2011.14443, mentioned for dimension-two embedded resolution. Titles and authors are not supplied in the record.
\end{itemize}
\end{document}
